% BHU PYQ -- Manually transcribed & error-corrected
% Paper: GRB-401: Regional Geography of Developed and Developing World (U.S.A. and China)
% Exam: B.Sc. (Hons.) Semester IV Examination, 2022-23

\documentclass[11pt,a4paper]{article}
\usepackage[a4paper,top=2.5cm,bottom=2.5cm,left=2.8cm,right=2.8cm,headheight=14pt]{geometry}
\usepackage{amsmath,amssymb}
\usepackage[T1]{fontenc}
\usepackage[utf8]{inputenc}
\usepackage{lmodern,microtype,fancyhdr,enumitem,hyperref,lastpage,parskip}

\hypersetup{colorlinks=true,urlcolor=black,linkcolor=black,
  pdftitle={GRB-401 Regional Geography -- BHU B.Sc. Sem IV 2022-23}}
\pagestyle{fancy}\fancyhf{}
\renewcommand{\headrulewidth}{0.4pt}\renewcommand{\footrulewidth}{0.4pt}
\fancyhead[L]{\small   \textit{Banaras Hindu University}}
\fancyhead[R]{\small   \textit{GRB-401: Regional Geography}}
\fancyfoot[C]{\small Page  \thepage\ of \pageref{LastPage}}


\newlist{parts}{enumerate}{2}
\setlist[parts,1]{label=(\alph*),leftmargin=*,itemsep=5pt,topsep=3pt}

\newcommand{\pts}[1]{\hfill   \textit{[#1]}}


\begin{document}

\begin{center}
  {\large\bfseries BANARAS HINDU UNIVERSITY}\\[2pt]
  {\normalsize B.Sc. (Hons.) --- Semester IV Examination, 2022-23}\\[6pt]
  \rule{0.8\linewidth}{0.5pt}\\[6pt]
  {\large\bfseries Paper No. GRB-401: Regional Geography of Developed and Developing World (U.S.A. and China)}\\[4pt]
  {\normalsize Time: 3 Hours \hfill Maximum Marks: 70}\\[2pt]
  {\small Ref. No.: 061364}
\end{center}
\rule{\linewidth}{0.4pt}
\smallskip



\noindent   \textit{Instructions: Answer   \textbf{five} questions in all. Question No. 1, carrying 20 marks, is compulsory. Remaining questions are of equal marks. Illustrate your answers with suitable maps wherever required.}
\bigskip




\noindent   \textbf{Question 1.} Write short answer of the following questions: \pts{20}

\begin{parts}
  \item Describe briefly the characteristics of developing countries. \pts{4}
  \item Give a brief geographical account of soil distribution in China. \pts{4}
  \item What factors are responsible for the development of the New-England Industrial Region of the USA? \pts{4}
  \item Describe briefly the distribution of iron-ore in China. \pts{4}
  \item Describe briefly the various bases used to classify developing and developed countries of the world. \pts{4}
\end{parts}

\medskip


\noindent   \textbf{Question 2.} Divide the USA into various physiographic regions and explain the characteristic features of the Interior Plains in detail. \pts{12.5}

\medskip


\noindent   \textbf{Question 3.} Give a geographical account of the petroleum resources of China. \pts{12.5}

\medskip


\noindent   \textbf{Question 4.} Give an account of the distribution of the Iron and Steel Industry in the USA. Describe the factors responsible for its concentration around the Great Lakes Region. \pts{12.5}

\medskip


\noindent   \textbf{Question 5.} Give a geographical account of population distribution in China. \pts{12.5}

\medskip


\noindent   \textbf{Question 6.} Enlist the major Agricultural Regions of the USA and describe the characteristic features of the Corn Belt in detail. \pts{12.5}

\medskip


\noindent   \textbf{Question 7.} Divide China into various geographical regions and explain any one of them in detail. \pts{12.5}

\medskip


\noindent   \textbf{Question 8.} Describe the various types of agricultural regions of China and their attributes. \pts{12.5}

\vfill\rule{\linewidth}{0.4pt}

\begin{center}\small
\href{https://bhu-exam.vercel.app/}{Click here to visit our website}\\[4pt]
\href{https://me-aryan.vercel.app/}{Click here to visit our contributor}\\[4pt]
\href{https://quashan-me.vercel.app/}{Click here to visit our contributor}
\end{center}
\end{document}
